\documentclass[10pt,a4paper]{amsart}
\usepackage[margin=19mm]{geometry}
\usepackage{amsmath,amssymb,mathtools}
\usepackage[scr=rsfs,cal=euler]{mathalfa}
\usepackage{xfrac}
\usepackage[unicode,hidelinks,bookmarks=false]{hyperref}
\newcommand{\ie}{i.e.\ }
\newcommand{\cf}{cf.\ }
\newcommand{\resp}{respectively\ }
\newcommand{\Subsection}{\subsection}
\providecommand{\nobreakdash}{\nobreak\hbox{-}}
\setlength{\emergencystretch}{2em}
\multlinegap0pt
\usepackage{bm}
\usepackage{bbm}
\usepackage{dsfont}
\usepackage{xspace}
\usepackage{cancel}
\usepackage{mathtools}
\usepackage{amsthm}
\makeatletter
\@ifundefined{thm}{\newtheorem{thm}{Thm}}{}
\@ifundefined{cor}{\newtheorem{cor}[thm]{Corollary}}{}
\@ifundefined{lem}{\newtheorem{lem}[thm]{Lemma}}{}
\makeatother
\newcommand{\bb}{\mathbb}
\newcommand{\ff}[1]{^{\underline{#1}}}
\newcommand{\mcal}{\mathcal}
\newcommand{\st}{\mathrm{c}}
\newcommand{\sym}{\mathrm{Sym}}
\title[Intersecting families of permutations with a fixed number of]{Intersecting families of permutations with a fixed number of cycles}
\author{Venkata Raghu Tej Pantangi}
\date{Source: arXiv:2608.20248v1 (CC BY 4.0)}
\begin{document}
\maketitle

\noindent\emph{Source and licence.} Intersecting families of permutations with a fixed number of cycles. Version arXiv:2608.20248v1 (CC BY 4.0), distributed under the Creative Commons Attribution 4.0 International licence, as recorded in the arXiv licence metadata for this exact version and shown on the abstract page https://arxiv.org/abs/2608.20248v1. Full rights notice: licenses/arxiv-2608.20248-CC-BY-4.0.txt.\par\smallskip

\noindent\textit{Editorial scope.} Counted unit: the Lem whose source label is \texttt{lem:HMlowerbound} in the article (source0.tex lines 520--523, source label \texttt{lem:HMlowerbound}), delivered together with the article's own complete original proof of it (source0.tex lines 524--581, one \texttt{proof} environment of 58 lines). No mathematics is written, repaired or re-derived: statement and proof are the article's own text line for line. Required context retained uncounted: lem:keyineq (lines 294--297); lem:keyineq1 (lines 324--327); eq:bonfterm (lines 468--470); eq:bonf (lines 475--477); cor:xjlongcycle (lines 503--509). Every definition, display and label of those lines is retained unchanged. Omitted from this package and not redistributed: 1--293, 298--323, 328--467, 471--474, 478--502, 510--519, 582--945. Nothing inside a retained excerpt is omitted or altered: every retained body is delivered exactly as the article's lines are written, so no omission receipt of any kind accompanies this package. Citations: the retained excerpts carry no citation command at all, so no bibliography entry of the article is transcribed and no reference is redistributed. Reference adaptation: none was needed and none was made; every reference inside the delivered unit resolves inside it, so no unresolved reference or citation is produced. Printed slips kept literal: no printed word, symbol, hypothesis or number of the counted statement or of its proof was altered. Layout and class: the article's own class is \texttt{amsart} with packages including amsmath, graphicx, hyperref, amsfonts, amssymb, amsthm, amscd, enumerate, todonotes; the wrapper is the lane's pinned \texttt{amsart} editorial wrapper at 10pt on A4, which restates the theorem-like environments the retained text needs and adds the article's own macro definitions that the retained text needs (\texttt{\textbackslash bb}, \texttt{\textbackslash ff}, \texttt{\textbackslash mcal}, \texttt{\textbackslash st}, \texttt{\textbackslash sym}), with extra wrapper packages (bm, bbm, dsfont, xspace, cancel, mathtools). The delivered numbering is the unit's own. Assets: no retained span contains a figure environment, a graphics-inclusion command, a PDF-page inclusion, a table or a TikZ picture; no image file is copied into this package, the package declares no asset dependency and no image file is redistributed with it. Licence: version arXiv:2608.20248v1 of this article is distributed under the Creative Commons Attribution 4.0 International licence, as recorded in the arXiv licence metadata for this exact version and shown on the abstract page https://arxiv.org/abs/2608.20248v1. A full rights notice is delivered at licenses/arxiv-2608.20248-CC-BY-4.0.txt. Characters: every retained excerpt is pure ASCII, so the delivered engine, XeLaTeX with its default Latin Modern text font, reports no missing character.
\begin{lem}\label{lem:keyineq}
If $\beta \in (0,1)$, $m,l,x$, $y$, $w$ are non-negative integers such that $m>(2e)^{e/(1-\beta)}$, $0<l\leq m^{\beta}$, $0\leq y\leq l$, $0\leq w\leq m^{\beta}$, and $0\leq x< m-w$ then 
\[ \st(m-w,l) \geq  (m-w-1)\ff{x} \left(\dfrac{(1-\beta)\ln m}{2m^{\beta}} \right)^{y} \st(m-w-x,l-y). \]
\end{lem}

\begin{lem}\label{lem:keyineq1}
If $\alpha \in(0, 1)$, $\delta \in (0,1/2)$ and $m,l,w,x$ are  integers such that\\ $m\geq \max\{(2e)^{e/(1-\alpha)}, e^{(1-\delta)/(\delta(1-\alpha)}\}$, $1\leq l \leq m^{\alpha}$, $0\leq x \leq l-1$, and $0\leq w \leq m^{\alpha}$, then 
\[\st(m-w,l) \leq \left(\dfrac{1}{1-\delta} \right)^{x} (m-w-1)\ff{x} \st(m-w-x,l).\]
\end{lem}

\begin{equation}\label{eq:bonfterm}
\bb{S}_{z}((r,s);\pi) = \sum\limits_{\substack{Z \subseteq V_{r,s}^{\pi}\\  |Z|=z }} \st(n-|Z|-1,k-\mu_{r,s}^{\pi}(Z))
\end{equation}

\begin{equation}\label{eq:bonf}
\sum\limits_{z=1}^{2m} (-1)^{z+1}\bb{S}_{z}((r,s);\pi) \leq |\sym(n,k)[(r,s);\pi]| \leq \sum\limits_{z=1}^{2m+1} (-1)^{z+1}\bb{S}_{z}((r,s);\pi),
\end{equation}

\begin{cor}\label{cor:xjlongcycle}
Let $n>k\geq1$, $\ell>1$ be integers. Further, let $r,s\in [n]$ and $\pi \in \sym(n,k)$ such that $(r,s,\pi)$ is admissible and the shortest cycle in $\pi$  has length $\ell$. Then the following hold:
\begin{enumerate}
\item If $r\neq s$ and $r,s$ lie in disjoint cycles of $\pi$, then $\mcal{X}_{J}$ has no cycles for all $J\subset V_{r,s}^{\pi}$ with $|J|<\ell$.
\item If $r=s$, then for all $J\subset V_{r,s}^{\pi}$ with $|J|<\ell$, the only cycle in $\mcal{X}_{J}$ is the loop on the vertex $r$.
\end{enumerate}
\end{cor}

\begin{lem}\label{lem:HMlowerbound}
Given $\xi\in (0,0.5)$ and $\alpha \in (0,1)$, for $n$ sufficiently large and $k\leq n^{\alpha}$, for each  $r,s\in [n]$, there exists $\pi \in \sym(n,k)$ such that $(r,s,\pi)$ is admissible and 
\[(1-1/e -\xi)\st(n-1,k-\delta_{r,s}) \leq |\sym(n,k)[(r,s);\pi]| \leq (1-1/e +\xi)\st(n-1,k-\delta_{r,s}).\] (Here $\delta_{r,s}$ is the Kronecker delta function.)
\end{lem}

\begin{proof}
We note that $n^{\alpha} n^{(1-\alpha)^{2}}=n ^{1-\alpha(1-\alpha)}<n$.
Therefore, since $k\leq n^{\alpha}$, there exists a permutation in $\sym(n,k)$ with all of its cycles of length at least $\lfloor n^{(1-\alpha)^{2}} \rfloor$. We shall refer to such a permutation as a {\it long-cycle permutation}.

Given $r,s\in[n]$, choose a long-cycle permutation $\pi$ as follows. If $r\neq s$ and $k\geq 2$, we pick $\pi$ such that $r$ and $s$ lie in distinct cycles of $\pi$. If $r\neq s$ and $k=1$, choose $\pi$ to be an $n$-cycle of the form $(r,a,s,\ldots)$, where $a\in[n]\setminus \{r,s\}$; such an $a$ exists for $n\geq3$. If $r=s$, choose any long-cycle permutation $\pi$.

Using \eqref{eq:bonf}, we find a lower bound on $|\sym(n,k)[(r,s);\pi]|$. 
Let $m$ be the largest integer such that $2m+1 < \min\{n^{\alpha}, \lfloor n^{(1-\alpha)^{2}}\rfloor \}$. Using \eqref{eq:bonf}, we have 
\begin{equation}\label{eq:bonfineq}
\sum\limits_{z=1}^{2m+1} (-1)^{z+1}\bb{S}_{z}((r,s);\pi)\geq |\sym(n,k)[(r,s);\pi]| \geq \sum\limits_{z=1}^{2m} (-1)^{z+1}\bb{S}_{z}((r,s);\pi),
\end{equation}
 where $\bb{S}_{z}((r,s);\pi)$ is as defined by \eqref{eq:bonfterm}. 

Consider $J \subset V_{r,s}^{\pi}=[n]\setminus \{r,\pi^{-1}(s)\}$ with $|J|\leq 2m+1$. 
We recall that the number $\mu_{r,s}^{\pi}(J)$ represents the number of cycles in $\mcal{X}_{J}$. Since $|J|\leq 2m+1 <\lfloor n^{(1-\alpha)^{2}}\rfloor$ and every cycle of $\pi$ has length at least $\lfloor n^{(1-\alpha)^{2}}\rfloor$, Corollary~\ref{cor:xjlongcycle} implies $\mu_{r,s}^{\pi}(J)=\delta_{r,s}$, where $\delta_{r,s}$ is the standard Kronecker delta function. Applying this in \eqref{eq:bonfterm}, we obtain

\begin{equation}\label{eq:bonftermlongpermint}
\bb{S}_{z}((r,s);\pi) = \binom{n-2}{z} \st(n-z-1,k-\delta_{r,s}),
\end{equation}
for all $1\leq z\leq 2m+1$.

Given $\epsilon\in (0,0.5)$ and $n$ sufficiently large, by Lemmas~\ref{lem:keyineq} and \ref{lem:keyineq1}, we have 
\[\dfrac{(1-\epsilon)^{z}}{(n-2)\ff{z}} \st(n-1,k-\delta_{r,s}) \leq \st(n-z-1,k-\delta_{r,s}) \leq \dfrac{1}{(n-2)\ff{z}} \st(n-1,k-\delta_{r,s}).\]
Applying the above along with \eqref{eq:bonftermlongpermint} yields
\begin{equation}\label{eq:bonftermlongperm}
\dfrac{(1-\epsilon)^{z}}{z!}\st(n-1,k-\delta_{r,s}) \leq \bb{S}_{z}((r,s);\pi) \leq \dfrac{1}{z!}\st(n-1,k-\delta_{r,s})
\end{equation}
for all $1\leq z\leq 2m+1$.
Applying the above in \eqref{eq:bonfineq}, we have
\begin{align*}
\dfrac{|\sym(n,k)[(r,s);\pi]|}{\st(n-1,k-\delta_{r,s})} &\geq \sum\limits_{i=1}^{m} \bb{S}_{2i-1}((r,s);\pi)-\sum\limits_{i=1}^{m} \bb{S}_{2i}((r,s);\pi) \\
&\geq \sum\limits_{i=1}^{m} \dfrac{(1-\epsilon)^{(2i-1)}}{(2i-1)!} - \sum\limits_{i=1}^{m} \dfrac{1}{(2i)!}  \\
&\geq 1+\sum\limits_{i=1}^{m} \dfrac{(1-\epsilon)^{(2i-1)}}{(2i-1)!} - \sum\limits_{i=0}^{m} \dfrac{1}{(2i)!}
\end{align*}
A similar argument leads to an upper bound, that is, we have
\begin{equation}\label{eq:partialsum0}
\begin{aligned}
1+\sum_{i=1}^{m}
\frac{(1-\epsilon)^{2i-1}}{(2i-1)!}
-\sum_{i=0}^{m}\frac{1}{(2i)!}
&\leq
\frac{|\sym(n,k)[(r,s);\pi]|}
     {\st(n-1,k-\delta_{r,s})} \\
&\leq
1+\sum_{i=1}^{m}
\frac{1}{(2i-1)!}
-\sum_{i=0}^{m}
\frac{(1-\epsilon)^{2i}}{(2i)!}.
\end{aligned}
\end{equation} for any given $\epsilon \in (0,0.5)$ and for all sufficiently large $n$. 

Set $A(x)=1+\sinh(1)-\cosh(1-x)$ and $B(x)=1+ \sinh(1-x)-\cosh(1)$. Fix $\xi\in (0,0.5)$, and by continuity of $A$ and $B$, choose $\epsilon \in (0,0.5)$ sufficiently small such that 
$A(\epsilon)>A(0)-\xi/2$  and $B(\epsilon)<B(0)+\xi/2$. We note that as $n\to \infty$, the upper bound in \eqref{eq:partialsum0} converges to $B(\epsilon)$ and the lower bound to $A(\epsilon)$. Therefore, for all $n$ sufficiently large, we have
\[B(0)+\xi>B(\epsilon)+\xi/2> \dfrac{|\sym(n,k)[(r,s);\pi]|}
     {\st(n-1,k-\delta_{r,s})}>A(\epsilon)-\xi/2>A(0)-\xi.\]

The result follows by $A(0)=1-1/e=B(0)$.
\end{proof}

\end{document}
